% =====================================================================
% THESIS DRAFT SECTIONS — generated 2026-06-04 for Overleaf integration
% =====================================================================

\section*{Methods: Spatial PDE Model}
\subsection{Spatial Reaction-Diffusion Model of Biofilm Depth Profiles}
\label{subsec:spatial_pde}

To connect the community-level ODE attractors to the three-dimensional FISH measurements,
we formulated a one-dimensional reaction-diffusion PDE in the depth coordinate $z$,
representing the axis perpendicular to the implant surface in the HOBIC biofilm system
\cite{heine2025}.
Let $\phi_i(z,t)$ denote the relative abundance (volume fraction) of guild $i$
at depth $z \in [0, L]$ and time $t$, where $L$ is the total biofilm thickness.
The model reads

\begin{equation}
  \frac{\partial \phi_i}{\partial t}
  = \underbrace{D_i \frac{\partial^2 \phi_i}{\partial z^2}}_{\text{self-diffusion}}
  + \underbrace{\sum_{j \neq i} \varepsilon_{ij}
      \frac{\partial^2 \phi_j}{\partial z^2}}_{\text{cross-diffusion}}
  + \underbrace{f_i(\boldsymbol{\phi})}_{\text{reaction}},
  \qquad i = 1, \ldots, 5,
\label{eq:pde_main}
\end{equation}

\noindent
where $D_i \geq 0$ is the self-diffusion coefficient of guild $i$,
$\varepsilon_{ij}$ is a cross-diffusion coupling constant (taken uniform,
$\varepsilon_{ij} = \varepsilon$, as an isotropic first-order approximation),
and $f_i(\boldsymbol{\phi})$ is the replicator-form reaction term inherited from the
fitted gLV attractor, evaluated pointwise at each depth level.
No-flux (Neumann) boundary conditions $\partial_z \phi_i = 0$ were imposed at both
$z = 0$ (substratum) and $z = L$ (biofilm--medium interface).
The system was initialised from a spatially uniform composition matching the
attractor fixed point of the respective community state (commensal-HOBIC, CH;
or dysbiotic-HOBIC, DH).

The diffusion coefficients $\mathbf{D} = (D_1, \ldots, D_5)$ and the scalar
cross-diffusion parameter $\varepsilon$ were estimated by minimising the sum of
squared residuals between the model steady-state depth profile $\hat{\phi}_i(z)$
and the experimental FISH z-profiles averaged across all fields of view:

\begin{equation}
  \mathcal{L}(\mathbf{D}, \varepsilon)
  = \sum_{i=1}^{5} \sum_{k=1}^{N_z}
    \Bigl[\hat{\phi}_i(z_k) - \phi_i^{\mathrm{FISH}}(z_k)\Bigr]^2,
\label{eq:loss}
\end{equation}

\noindent
where $z_k$ are the $N_z$ evenly spaced depth bins into which the confocal stack
was resampled.
The optimisation was performed with a bounded L-BFGS-B solver;
$D_i$ was constrained to $[10^{-5},\, 0.5]$ in dimensionless solver units
and $\varepsilon \in \{10^{-4}, 3\times10^{-4}, 10^{-3}\}$ was treated as a
fixed hyperparameter swept across conditions.
Grid-resolution convergence was assessed over $N_z \in \{24, 30, 32, 40\}$.

For the CH attractor the model was evaluated in four configurations:
baseline (self-diffusion only, $\varepsilon = 0$) and three cross-diffusion
variants ($\varepsilon = 10^{-4}$, $3\times10^{-4}$, $10^{-3}$) at each grid
resolution.
The best CH result was obtained at $N_z = 32$, $\varepsilon = 3\times10^{-4}$,
yielding a fitting loss of $\mathcal{L} = 0.0145$, a 52\% reduction relative to
the baseline ($\mathcal{L} = 0.0300$), demonstrating that cross-diffusion is
essential for reproducing the CH depth structure.
The dominant mobile species in this configuration was \textit{A.~naeslundii}
($D = 0.500$ solver units; note that this value coincides with the upper bound of the
search domain, indicating $D_{An} \geq 0.5$~\todo{expand search range and re-run before
reporting a calibrated physical value}) and \textit{S.~oralis}
($D = 3.75 \times 10^{-3}$), consistent with their roles as early colonisers that
penetrate the outer biofilm layers \cite{kolenbrander2010}.
\textit{V.~dispar}, \textit{F.~nucleatum}, and \textit{P.~gingivalis} were effectively
immobile ($D \approx 10^{-5}$), in accord with the deeper, anaerobic niche occupied
by late colonisers \cite{haffajee2008}.
Dimensionless solver units have not yet been converted to physical units (mm$^2$\,h$^{-1}$);
quantitative comparison with literature diffusivity values requires this calibration
step \todo{verify physical unit conversion factor}.

The DH attractor presented a qualitatively different behaviour: the best fitting loss
across all successful configurations was $0.091$--$0.103$, approximately $7\times$ higher
than the best CH loss, and runs at $N_z \geq 32$ failed to converge for the majority of
parameter initialisations.
Crucially, adding cross-diffusion did not improve the DH fit
($\mathcal{L}_{\mathrm{xdiff}} \geq 0.103$ vs.\ $\mathcal{L}_{\mathrm{baseline}} = 0.091$).
This outcome is interpreted not as a model failure but as a positive finding:
the DH state corresponds to the confluent-mat morphology identified in the FISH data,
where the lateral coefficient of variation is $\mathrm{CV} \approx 0.33$ and
depth gradients are weak.
In this regime reaction kinetics dominate over transport, rendering the diffusion
coefficients non-identifiable from the observed profiles --- a spatial manifestation
of the community restructuring that homogenises the biofilm under dysbiotic conditions.

\section*{Results: Diffusion Coefficient Fitting}
\subsection{Diffusion Coefficient Fitting}
\label{subsec:diffusion_fit}

The one-dimensional reaction-diffusion PDE was fitted to depth-resolved species
profiles for the commensal-HOBIC (CH) and dysbiotic-HOBIC (DH) attractors by
minimising a mean-squared loss between simulated and observed $z$-profiles.
A sweep over grid resolution ($N_z \in \{24, 30, 32, 40, 48\}$), finite-difference
step ($\varepsilon \in \{10^{-3}, 3\times10^{-4}, 3\times10^{-3}\}$), and the
presence or absence of a cross-diffusion term was conducted to assess numerical
sensitivity and structural identifiability (Table~S1).

\paragraph{CH attractor: well-determined sparse diffusion structure.}
Without cross-diffusion, all CH configurations converge to a loss of approximately
\textbf{0.030} regardless of $\varepsilon$, with \textit{F.~nucleatum} and
\textit{Veillonella}/\textit{Prevotella} (Vd/Vp) as the dominant mobile species
($D \approx 6$--$7\times10^{-3}$) and the remaining three species effectively
immobile ($D \sim 10^{-5}$).
Introducing the cross-diffusion coupling at $N_z = 32$, $\varepsilon = 3\times10^{-4}$
cuts the loss to \textbf{0.0145} — a \textbf{52\%} reduction relative to the
baseline (Fig.~\ref{fig:fit_CH}).
Under this best-fit configuration the mobile-species assignment reverses:
\textit{A.~naeslundii} becomes the dominant diffusing species
($D_{\mathrm{An}} = \mathbf{0.500}$, solver units), with \textit{S.~oralis}
at $D_{\mathrm{So}} = 3.75\times10^{-3}$, while Vd/Vp, \textit{F.~nucleatum},
and \textit{P.~gingivalis} are effectively immobile ($D \sim 10^{-5}$).
The advection velocity is small ($u = 1.76\times10^{-4}$), confirming that the
improvement stems from cross-diffusive coupling rather than bulk transport.
A note of caution applies: $D_{\mathrm{An}} = 0.500$ is at the upper boundary of
the search box, indicating that the optimiser's box constraint is active and the
true optimal value may be higher.
This figure should therefore be read as a lower bound ($D_{\mathrm{An}} \geq 0.5$,
solver units); quantitative conversion to physical units (mm$^2$ h$^{-1}$)
requires expanding the search range in a follow-up run.

Grid resolution has a monotonic but moderate effect on CH performance: coarsening
from $N_z = 32$ to $N_z = 24$ raises the (no-xdiff) loss from $0.030$ to $0.038$,
while $N_z = 48$ diverges or fails convergence entirely, placing the sweet spot
firmly at $N_z = 32$.

\paragraph{DH attractor: reaction-dominated regime.}
The DH sweep yields a qualitatively different picture (Fig.~\ref{fig:fit_DH}).
The best successful run (\texttt{xdiff\_nz40\_eps1e-3}) achieves a loss of
\textbf{0.103}, and the floor across all successful DH configurations lies in the
range \textbf{0.091--0.103} — approximately \textbf{7$\times$ higher} than the
best CH loss.
Crucially, enabling cross-diffusion does \emph{not} improve DH performance:
the xdiff DH runs (0.103--0.103) match or slightly exceed the best non-xdiff DH
loss (0.099, though that run did not converge), and the DH loss surface appears
flat across all $\varepsilon$ settings.
At $N_z \geq 32$ the majority of DH runs fail to converge; failure rates are
higher than for any CH configuration, and $N_z = 48$ is consistently the worst
condition (loss $\approx 0.154$, all three runs failed or degraded).
The fitted $D$ values in successful DH runs show \textit{S.~oralis} and
\textit{A.~naeslundii} co-mobile at similar levels
($D_{\mathrm{So}} \approx D_{\mathrm{An}} \approx 0.026$--$0.041$), with
\textit{P.~gingivalis} secondarily mobile ($D_{\mathrm{Pg}} \approx 0.005$--$0.010$)
and Vd/Vp immobile — but the variance across converged runs is large, reflecting
poor identifiability.

This high loss floor and convergence fragility are interpreted as a scientifically
affirmative finding rather than a model failure.
FISH imaging of the DH state (Chapter~\ref{ch:fish}) shows a confluent mat
morphology with low lateral heterogeneity (coefficient of variation $\approx 0.33$),
indicating that the dysbiotic community is spatially homogeneous along the biofilm
depth axis.
When spatial structure is dominated by reaction kinetics rather than transport,
diffusion coefficients become non-identifiable from $z$-profile data alone — the
reaction-diffusion PDE has no lever to distinguish species by position.
The 7$\times$ gap in loss between CH and DH is therefore consistent with the
interpretation that the commensal-HOBIC state organises species spatially through
differential mobility, whereas the dysbiotic-HOBIC state collapses this
organisation, producing a reaction-dominated regime in which fitting diffusion
coefficients is inherently underdetermined.

\section*{Results: FISH 3D Co-occurrence Analysis}
\subsection{FISH 3D Co-occurrence Analysis}
\label{sec:fish_cooccurrence}

To obtain an independent spatial read of inter-species relationships in
HOBIC biofilms, we quantified pairwise spatial co-occurrence from
three-dimensional FISH image stacks using Pearson correlation of
species-resolved voxel intensity profiles (Figure~\ref{fig:fish_cooccurrence}).
The analysis covered $n = 137$ field-of-view observations in the commensal-HOBIC
(CH) attractor and $n = 157$ in the dysbiotic-HOBIC (DH) attractor.

Across all ten species pairs, raw FISH co-occurrence was weak: the majority
of mean Pearson $r$ values fell below $|r| < 0.20$ in both attractors.
No pair showed strong positive spatial association, indicating that the
five focal species (\textit{S.~oralis}, \textit{A.~naeslundii},
\textit{Vd/Vp}, \textit{F.~nucleatum}, \textit{P.~gingivalis}) do not
form tight spatial clusters but are instead distributed across the biofilm
volume without pronounced proximity preferences.
When compared against the sign of the symmetric interaction matrix
$A_\mathrm{sym}$ inferred from the Hamilton replicator model,
sign concordance between FISH $\mathrm{sign}(\bar{r})$ and
$\mathrm{sign}(A_\mathrm{sym})$ was 5 of 10 pairs in CH and 6 of 10 pairs
in DH.

The most notable discordance involves the \textit{F.~nucleatum}--\textit{P.~gingivalis}
(Fn-Pg) pair.  The model assigns the largest interaction magnitude in the
entire $A_\mathrm{sym}$ matrix to this pair ($A_\mathrm{sym} = +1.38$ in CH,
$+9.73$ in DH), predicting a cooperative or co-occurring relationship,
yet FISH spatial correlation is near zero and negative in both attractors
($\bar{r} = -0.046$ in CH; $\bar{r} = -0.126$ in DH).
This apparent contradiction is mechanistically interpretable: \textit{F.~nucleatum}
is established in the literature as a ``bridging species'' that facilitates
\textit{P.~gingivalis} colonisation through adhesin--fimbriae binding and
metabolic cross-feeding, processes that operate through transient or
sequential contact rather than persistent spatial co-localisation.
The ecological interaction strength captured by $A_\mathrm{sym}$ therefore
reflects a temporal or chemical dependency rather than a static proximity
relationship, and the absence of a FISH co-occurrence signal is consistent
with, not contradictory to, the model.

A second notable case is the \textit{S.~oralis}--\textit{P.~gingivalis}
(So-Pg) pair.  The model undergoes a sign reversal between attractors
($A_\mathrm{sym} = -0.33$ in CH, $+0.85$ in DH), whereas FISH shows
a consistently negative correlation in both conditions
($\bar{r} \approx -0.20$ in CH, $-0.28$ in DH).  The FISH data thus do not
support the attractor-specific sign flip predicted by the dynamical model,
and this pair is flagged as a candidate for further investigation.
In contrast, the \textit{S.~oralis}--\textit{Vd/Vp} pair shows one of
the cleaner agreements: $A_\mathrm{sym} = +0.86$ (CH) and
$\bar{r} = +0.15$, suggesting that early-coloniser spatial proximity
is partially consistent with the inferred cooperative interaction.

Temporal dynamics of the depth-resolved FISH profiles are summarised in
Figure~\ref{fig:fish_zprofile_temporal}.
The $z$-profile structure is discussed in detail in the context of the
reaction-diffusion PDE analysis (Section~\ref{sec:spatial_pde});
briefly, DH biofilms exhibit lateral coefficient of variation
$\mathrm{CV} \approx 0.33$, indicative of a confluent, spatially
homogenised mat, whereas CH biofilms retain greater lateral heterogeneity,
consistent with a sparser, spatially structured community.

Taken together, these results indicate that FISH spatial co-occurrence
is a weak proxy for ecological interaction strength in HOBIC biofilms.
The modest concordance fraction (5--6 of 10 pairs) is expected given
that the gLV interaction matrix encodes population-level dynamics
integrated over time, while FISH captures a single spatial snapshot.
The analysis nonetheless provides a useful independent constraint: the
pairs that do agree (An-Pg, Vd/Vp-Pg, So-Vd/Vp) offer cross-modal
corroboration, while the Fn-Pg discordance sharpens the biological
interpretation of the largest inferred interaction in the model.

\section*{Results: LOO Stability + Discussion}
\subsection{LOO Cross-Validation Stability and Prior Asymmetry}

To assess whether the fitted interaction matrix $\mathbf{A}$ is uniquely determined by the Dieckow longitudinal data or is artificially constrained by the metabolic sign prior, we performed a systematic stability analysis across all ten leave-one-out (LOO) folds. For each off-diagonal entry $A_{ij}$, we classified its behaviour into four categories: \emph{aligned} (data and prior agree on sign), \emph{data-driven} (no prior penalty active, yet sign is stable across folds), \emph{prior-constrained} (prior overrides data), and \emph{uncertain} (sign inconsistent across folds).

Of the 90 off-diagonal entries in the $10 \times 10$ interaction matrix, 52 (57.8\%) were classified as aligned, 36 (40.0\%) as data-driven, 1 (1.1\%) as prior-constrained, and 1 (1.1\%) as uncertain. The consensus sign matrix was unanimous (10/10 LOO folds) for all but two entries: $A_{[1,2]}$ achieved 8/10 consensus and $A_{[7,1]}$ achieved 7/10 consensus, indicating near-complete sign identifiability from the data alone. Coefficient of variation across folds was low: for representative strong interactions, $A_{[1,3]}$ had mean $= 1.757$ and standard deviation $= 0.113$ (CV $= 0.064$), and $A_{[0,1]}$ had mean $= 1.588$ and standard deviation $= 0.058$ (CV $= 0.037$). These values are consistent across all LOO configurations, confirming that the interaction matrix is not overfit to any single patient trajectory.

The prior asymmetry analysis revealed that the metabolic sign prior acts predominantly in reinforcement rather than in correction: 94.4\% of entries for which a prior was active also had independent data support for the same sign (aligned category). Only a single entry was classified as prior-constrained, indicating that the L1/L2/L3 metabolic hierarchy imposes negligible distortion on the posterior interaction structure. Entries without a prior penalty (data-driven, 40\%) nonetheless achieved stable sign consensus, demonstrating that the ecological data are sufficiently informative to determine interaction sign independently of the AGORA-derived metabolic constraints.

\subsection{Integrative Interpretation of Attractor Dynamics}

Taken together, the LOO stability, FISH co-occurrence, and spatial PDE results constitute a multilevel validation of the fitted interaction model. The LOO analysis establishes that the $\mathbf{A}$-matrix sign structure is robustly identified from ten-patient longitudinal data: near-unanimous fold consensus and CV $\ll 0.1$ for most entries directly rebut the underdetermination concern inherent in fitting $\mathcal{O}(100)$ parameters from short microbial time series. The spatial PDE results provide an orthogonal physical constraint. The commensal-HOBIC (CH) attractor yields a well-determined sparse diffusion structure, with cross-diffusion reducing fitting loss by 52\% relative to the diagonal baseline; \textit{A.\ naeslundii} and \textit{S.\ oralis} emerge as the mobile early-coloniser species, consistent with their established roles in initial biofilm formation. By contrast, the dysbiotic-HOBIC (DH) attractor resists diffusion fitting (loss floor $\approx 7\times$ higher, majority of high-resolution runs failing convergence), which is not a modelling failure but an affirmative finding: FISH imaging shows that DH is characterised by a confluent mat morphology with lateral coefficient of variation $\approx 0.33$, indicating that reaction kinetics rather than species transport govern spatial organisation. The PDE thus confirms at the physical level what FISH observes morphologically --- dysbiosis collapses spatial structure into a reaction-dominated regime in which diffusion coefficients become non-identifiable. FISH spatial co-occurrence shows modest overall sign agreement with the symmetric interaction matrix $\mathbf{A}_\mathrm{sym}$ (5--6 of 10 pairs per attractor), and the largest discordance --- \textit{Fn}--\textit{Pg}, where the strongest model coupling ($A_\mathrm{sym} = +9.73$ in DH) coincides with near-zero or negative spatial co-occurrence --- is mechanistically interpretable: \textit{F.\ nucleatum} facilitates \textit{P.\ gingivalis} colonisation through temporal bridging and chemical signalling rather than physical co-localisation, so spatial proximity is not expected to reflect ecological interaction sign for this pair. The three analytical pillars therefore converge on a coherent narrative: metabolic sign priors constrain but do not dominate the ODE inference; LOO stability confirms the inferred network is reproducible across patient subsets; and the spatial data confirm that the transition from commensal to dysbiotic attractor is accompanied by a physical collapse from diffusion-structured to reaction-dominated biofilm organisation.
